% Generado por scripts/validar_modelo_financiero.py — no editar a mano.
\begin{table}[t]
  \centering
  \footnotesize\setlength{\tabcolsep}{3.5pt}
  \caption{Flujo de caja del escenario pesimista, en USD. El flujo neto resta de los ingresos los costos, el impuesto a las ganancias y la inversión en capital de trabajo; en el año 5 esa inversión es negativa porque se recupera. Fuente: elaboración propia sobre la planilla del modelo financiero.}
  \label{tab:financiero-pesimista-flujo}
  \begin{tabular}{c r r r r r r r}
    \hline
    \textbf{Año} & \textbf{Ingresos} & \textbf{Costos} & \textbf{Impuesto} & \textbf{Cap. de trabajo} & \textbf{Flujo neto} & \textbf{Descontado} & \textbf{Acumulado} \\
    \hline
    0 & 0,00 & 1.715,56 & 0,00 & 0,00 & (1.715,56) & (1.715,56) & (1.715,56) \\
    1 & 0,00 & 33.948,91 & 0,00 & 0,00 & (33.948,91) & (27.159,12) & (35.664,47) \\
    2 & 0,00 & 39.239,62 & 0,00 & 0,00 & (39.239,62) & (25.113,36) & (74.904,09) \\
    3 & 35.200,00 & 64.670,30 & 0,00 & 5.866,67 & (35.336,97) & (18.092,53) & (110.241,06) \\
    4 & 62.400,00 & 98.382,01 & 0,00 & 4.533,33 & (40.515,35) & (16.595,09) & (150.756,41) \\
    5 & 86.800,00 & 113.943,86 & 0,00 & (10.400,00) & (16.743,86) & (5.486,63) & (167.500,27) \\
    \hline
  \end{tabular}
\end{table}
